\chapter{Conclusions}

\section{Summary of Key Findings}

This thesis presented a comprehensive methodology for quantifying spatial openness in residential 
properties through the integration of computer vision techniques and large-scale data analysis. By 
analyzing 4,004 rental properties across Tokyo's 23 wards, we developed a multifaceted evaluation 
approach that combines 2D visibility graph analysis (VGA) from floor plans with 3D architectural 
element extraction from interior photographs.

Our analysis revealed several significant findings regarding the evolution and distribution of 
spatial openness in Tokyo's residential architecture:

\subsection{Temporal Evolution of Spatial Design}

The temporal analysis uncovered a clear paradigm shift in residential design philosophy between the 
1960s and 1970s. Properties constructed in the 1960s consistently demonstrated higher VGA values 
across all metrics, with the most pronounced differences observed in the standard deviation of 
spatial connectivity (effect size $\eta^2 = 0.0609$, medium effect). This suggests that post-war 
reconstruction era housing featured more spatially integrated and varied layouts compared to the 
standardized designs that emerged from the 1970s onward.

The progressive compartmentalization of living spaces continued through the 2000s, particularly 
evident in living room designs where mean VGA values declined by up to 11.72\% from the 1960s to 
2000s. However, this trend stabilized by the 1990s, indicating that modern residential design has 
reached an equilibrium between openness and functional separation. Notably, the 2010s marked a 
partial reversal of this trend through increased window ratios (up to 30.54\% higher than previous 
decades), suggesting a renewed emphasis on natural light and visual connection to the exterior 
environment.

\subsection{Geographic Patterns of Openness}

The geospatial analysis revealed that spatial openness characteristics are not uniformly distributed 
across Tokyo but rather follow distinct patterns influenced by urban development and housing 
typology. Central Tokyo areas demonstrated higher window ratios, likely driven by high-rise 
apartment developments that prioritize views and natural light. In contrast, suburban areas showed 
lower VGA values, reflecting traditional Japanese housing preferences for privacy and clearly 
delineated functional spaces.

The scattered distribution of high VGA properties throughout the city suggests that open floor plans 
are not concentrated in specific districts but rather represent diverse architectural responses to 
local conditions, development timing, and target demographics. This finding challenges assumptions 
about uniform architectural trends within urban zones.

\subsection{Independence of 2D and 3D Openness Measures}

A critical finding emerged from the correlation analysis: 2D spatial connectivity (VGA) and 3D 
architectural elements (window, ceiling, floor, and wall ratios) operate as independent dimensions 
of openness. The lack of correlation between these measures indicates that horizontal spatial flow 
and vertical spatial qualities can be optimized separately, offering designers multiple pathways to 
create a sense of openness.

This independence was further validated through visual examples showing properties with high VGA but 
low window ratios achieving openness through layout design, while others with compartmentalized 
floor plans compensated through extensive glazing. This finding has important implications for 
design flexibility in constrained urban environments.

\subsection{Market Valuation of Spatial Openness}

The regression analysis demonstrated that spatial openness characteristics contribute moderately to 
rental price determination ($R^2 = 0.3$). Among the openness indicators, whole-property VGA 
statistics showed the highest predictive importance, suggesting that overall spatial connectivity 
matters more to market valuation than room-specific features. The retention of all whole-property 
VGA indicators (minimum, maximum, mean, and standard deviation) after multicollinearity reduction 
further emphasizes their distinct contributions to perceived value.

However, the disconnect between objective openness measures and subjective impression scores 
highlights the complexity of spatial perception. While our quantitative indicators capture important 
spatial characteristics, factors such as interior design, furniture arrangement, and aesthetic 
elements play equally important roles in shaping user experiences and preferences.

\section{Future Work}

While this study advances the quantification and analysis of spatial openness in residential 
environments, several limitations and areas for improvement have been identified that warrant future 
investigation:

\subsection{VGA Methodology Refinement}

The current VGA implementation faces accuracy limitations due to the granularity constraints of 
grid-based discretization. Our approach, while scalable, does not achieve the precision of 
ray-casting visibility algorithms used in specialized architectural software. Future work should 
explore:

\begin{itemize}
    \item Development of hybrid approaches that balance computational efficiency with accuracy
    \item Alignment with established architectural analysis benchmarks to validate results
    \item Investigation of adaptive grid resolutions that increase density in areas of high spatial 
    complexity
    \item Integration of more sophisticated visibility algorithms that account for partial occlusions 
    and graduated visual connections
\end{itemize}

\subsection{3D Indicator Robustness}

The extraction of 3D openness indicators from interior photographs revealed sensitivity to camera 
angles and shooting positions. Properties photographed from extreme high or low angles showed biased 
ratios for floor and ceiling elements, potentially leading to misrepresentative conclusions. Future 
research should address:

\begin{itemize}
    \item Development of angle-correction algorithms to normalize architectural element ratios
    \item Direct extraction of lighting conditions using specialized computer vision models
    \item Integration of depth estimation techniques to better understand spatial volumes
    \item Creation of standardized photography protocols for rental listings to ensure consistency
\end{itemize}

\subsection{Causal Analysis Framework}

The current study focused on correlational relationships between openness indicators and various 
outcomes. However, correlation does not imply causation, and establishing causal links would provide 
more actionable insights for design and policy decisions. Future work should investigate:

\begin{itemize}
    \item Development of causal inference frameworks to understand the directional relationships 
    between spatial design choices and outcomes
    \item Longitudinal studies tracking how renovations affecting openness impact rental prices and 
    tenant satisfaction
    \item Natural experiments leveraging regulatory changes or development patterns
    \item Integration of additional confounding variables such as neighborhood characteristics and 
    building amenities
\end{itemize}

The moderate $R^2$ values in our regression models suggest that while openness indicators contribute 
to rental pricing, many other factors remain unaccounted for. A more comprehensive causal framework 
would help isolate the specific impact of spatial openness from other property attributes.

\subsection{Scale and Generalizability}

The dataset of 4,004 properties, while substantial for a study involving detailed spatial analysis, 
represents a limited sample given Tokyo's vast rental market. The filtering process necessary to 
ensure data quality further constrained the sample size. Future research should address:

\begin{itemize}
    \item Expansion to larger datasets through improved automation of data processing pipelines
    \item Validation of findings across different cities and cultural contexts
    \item Development of transfer learning approaches to apply models trained on high-quality data to 
    lower-quality datasets
    \item Investigation of temporal stability through repeated sampling over multiple years
\end{itemize}

% \subsection{Integration of Subjective Experiences}

% The weak correlation between objective openness measures and impression evaluation scores highlights 
% the need for better integration of subjective spatial experiences. Future work should explore:

% \begin{itemize}
%     \item Development of multi-modal evaluation frameworks that combine objective measures with user 
%     perception data
%     \item Virtual reality studies to understand how different openness characteristics impact spatial 
%     experience
%     \item Ethnographic research to capture cultural and demographic variations in openness preferences
%     \item Machine learning approaches that can predict subjective responses from combinations of 
%     objective features
% \end{itemize}

% \subsection{Methodological Innovations}

% Several methodological improvements could enhance the robustness and applicability of this research:

% \begin{itemize}
%     \item Development of uncertainty quantification methods for both VGA and semantic segmentation 
%     outputs
%     \item Creation of ensemble approaches that combine multiple computer vision models for more 
%     reliable architectural element detection
%     \item Investigation of graph neural networks for more sophisticated spatial relationship modeling
%     \item Integration of natural language processing to analyze property descriptions and tenant 
%     reviews
% \end{itemize}

These future directions would not only address current limitations but also expand the scope and 
impact of spatial openness research, ultimately contributing to better-designed living environments.
